\bookchapter{Capitolo 22: Come funziona l'IA nel telefono}

\chapterintro{Benvenuti alla “sparizione magica” più cool del telefono, il protagonista è l'intelligenza artificiale. Le foto “bell'e pronte” che usi ogni giorno, l'assistente vocale pronto a rispondere, la tastiera che più la usi più ti capisce: dietro c'è questo cervello invisibile.}

\booksection{Fotocamera AI: la tua “camera oscura di post-produzione” si è trasferita nel telefono}

L'odierna fotografia computazionale AI assomiglia più a un team professionale di ritoccatori e fotografi stipato direttamente nel tasto otturatore. Prima che tu prema l'otturatore, l'IA ha già cominciato a lavorare: scansiona la scena e riconosce “questo è cielo azzurro”, “questa è una faccia”, “questa è una scena notturna”.

\booksection{Assistente vocale: come fa il tuo “segretario vocale” a capirti?}

\begin{enumerate}
  \item \textbf{Sentire chiaro (riconoscimento vocale)}: trasforma le onde sonore in testo.
  \item \textbf{Capire (comprensione del linguaggio naturale)}: comprende l'intento ed estrae le entità chiave.
  \item \textbf{Eseguire e rispondere}: richiama l'app o il servizio corrispondente e poi tramite sintesi vocale te lo "dice".
\end{enumerate}

\booksection{IA on-device vs. IA cloud}

L'IA on-device mantiene il calcolo nel telefono: è veloce, rispetta la privacy e funziona offline. L'IA cloud richiede connessione, ma può gestire compiti più complessi.

\booksection{Grandi modelli IA e telefono: il tuo assistente intelligente personale}

Nel 2025-2026, l'esplosione dei grandi modelli IA (come DeepSeek R1, ChatGPT-5, Wenxin Yiyan 5.0, Tongyi Qianwen 3, ecc.) ha cambiato profondamente il modo in cui si usa il telefono. Il telefono non è più solo uno strumento che esegue comandi, ma un partner intelligente in grado di comprendere, ragionare e creare. I chip IA on-device (con potenza di calcolo NPU oltre gli 80 TOPS) hanno permesso per la prima volta di eseguire l'inferenza dei grandi modelli localmente sul telefono, usando l'IA anche senza connessione.

\booksubsection{Che cos'è un grande modello IA?}

Un grande modello IA (Large Language Model, LLM) è un sistema IA addestrato con enormi quantità di dati testuali, in grado di comprendere e generare linguaggio naturale. Puoi immaginarlo come un “super studente modello” che ha letto quasi tutti i testi su Internet, capace di rispondere a domande, scrivere articoli, tradurre, programmare e analizzare dati.

\booksubsection{Applicazioni dei grandi modelli IA sul telefono:}

\begin{itemize}
  \item \textbf{Assistenti conversazionali}: App come DeepSeek, ChatGPT, Doubao, Kimi, con cui puoi fare domande, tradurre, scrivere testi e riassumere articoli in qualsiasi momento. Più intelligenti dei motori di ricerca tradizionali — non ti danno solo link, ma direttamente la risposta.
  \item \textbf{Scrittura e creazione}: lascia che l'IA ti aiuti a scrivere email, riassumere verbali di riunione, rifinire articoli, generare testi creativi.
  \item \textbf{Generazione e modifica di immagini}: l'IA elimina con un clic i passanti dalle foto, espande i bordi dell'immagine, converte lo stile o genera immagini a partire da descrizioni testuali.
  \item \textbf{Traduzione in tempo reale}: traduzione in tempo reale delle chiamate, traduzione di dialoghi faccia a faccia, traduzione di foto: niente più barriere linguistiche in viaggio all'estero.
  \item \textbf{Supporto allo studio}: l'IA può spiegare concetti, proporre esercizi, correggere i compiti: è un ripetitore personale sempre con te.
\end{itemize}

\booksubsection{Grandi modelli on-device: il telefono fa girare l'IA da solo}

I più recenti telefoni di punta possono già eseguire localmente piccoli grandi modelli IA (in genere circa 7B parametri). I vantaggi sono:

\begin{itemize}
  \item \textbf{Non serve rete}: le funzioni IA funzionano anche senza connessione, in aereo o in metro.
  \item \textbf{Privacy più sicura}: le tue conversazioni e i tuoi dati non vengono caricati sul cloud, ma elaborati interamente sul telefono.
  \item \textbf{Risposta più veloce}: senza latenza di rete, le risposte dell'IA sono quasi istantanee.
  \item \textbf{Limiti attuali}: le capacità dei modelli on-device non eguagliano ancora quelle dei grandi modelli cloud; i compiti complessi richiedono ancora la connessione. Ma con l'aumento della potenza di calcolo IA dei chip, i modelli on-device diventeranno sempre più forti.
\end{itemize}

\booksubsection{Integrazione IA a livello di sistema:}

\begin{itemize}
  \item \textbf{Apple Intelligence}: sistema IA introdotto in iOS 18, integra strumenti di scrittura, bacchetta magica per immagini, Siri potenziato, riassunto delle notifiche, ecc. Priorità all'elaborazione on-device, con cloud solo per i compiti complessi.
  \item \textbf{Modello Pangu di Huawei}: integrato in HarmonyOS, l'assistente Xiaoyi ottiene le capacità dei grandi modelli, supporta riassunti di documenti, riconoscimento intelligente dello schermo, fotoritocco IA, ecc.
  \item \textbf{Xiaomi Surge AI}: capacità IA a livello di sistema, supporta ricerca immagini IA, riassunti IA, sottotitoli in tempo reale IA, dialogo con grande modello on-device, ecc.
  \item \textbf{Samsung Galaxy AI}: traduzione in tempo reale delle chiamate, riassunti IA delle note, generazione e modifica di immagini, ricerca circolare, ecc.
\end{itemize}

\booksection{Alcuni scenari applicativi tipici e operazioni pratiche}

Per farti percepire in modo più intuitivo i “superpoteri” dell'IA sul telefono, ecco cinque scenari concreti che mostrano come l'IA “accorre in soccorso” nella vita quotidiana:

\booksubsection{Scenario 1: traduzione in tempo reale — l'interprete personale per i viaggi all'estero}

\textbf{Prestazioni reali}: a Parigi non sai il francese al momento di ordinare? L'app “Traduzione” di iPhone supporta la modalità dialogo faccia a faccia tra due persone, ciascuno vede i sottotitoli nella propria lingua, ed è possibile persino la traduzione simultanea. Sul fronte Android, la traduzione in tempo reale delle chiamate Samsung Galaxy AI ti permette di parlare in cinese al telefono mentre l'interlocutore ascolta in inglese, e viceversa.

\textbf{Principio alla base}: riconoscimento vocale (ASR) traduzione automatica (NMT) sintesi vocale (TTS), una pipeline in tre fasi con latenza end-to-end ormai contenibile sotto 1 secondo.

\booksubsection{Scenario 2: eliminazione IA — salvare le foto turistiche “rovinate” dai passanti}

\textbf{Prestazioni reali}: hai scattato un selfie stupendo in un luogo di interesse, ma lo sfondo è pieno di altri turisti? Con “Eliminazione magica” seleziona i passanti e il telefono li riconosce e li “cancella”, riempiendo intelligentemente lo sfondo. I risultati sono spesso sorprendenti.

\textbf{Principio alla base}: il modello di inpainting prima identifica ed elimina l'area selezionata, poi “immagina” un contenuto ragionevole sulla base dei pixel circostanti e della comprensione semantica. Il modello IA alla base è addestrato su milioni di immagini.

\booksubsection{Scenario 3: fotoritocco IA — diventa un “maestro della fotografia” con un clic}

\textbf{Prestazioni reali}: hai scattato una foto panoramica grigia in una giornata nuvolosa? Il “Cielo magico” IA può sostituire il cielo con un cielo sereno, un tramonto o persino un cielo stellato. L'espansione IA può trasformare una composizione verticale in orizzontale, “immaginando” le scene fuori dall'inquadratura.

\textbf{Principio alla base}: modello di diffusione (Diffusion Model) basato sull'IA generativa (Generative AI), che genera nuovi pixel dopo aver compreso il contenuto semantico dell'immagine. È la “frontiera più avanzata” dell'elaborazione di immagini IA attuale.

\booksubsection{Scenario 4: elaborazione documenti — il moltiplicatore di efficienza per chi lavora in ufficio}

\textbf{Prestazioni reali}: finita la riunione, passa la registrazione all'IA e otterrai automaticamente un verbale strutturato, con estrazione delle azioni da fare e generazione di un riassunto. Un report PDF inglese di 10 pagine può essere sintetizzato nei suoi punti chiave in 10 secondi.

\textbf{Principio alla base}: i grandi modelli a contesto lungo (Long Context LLM) possono “leggere” documenti di decine di migliaia di parole, effettuando ragionamento multi-turno ed estrazione di informazioni. Apple Intelligence, Huawei Xiaoyi, WPS AI, Tencent Yuanbao, ecc. hanno questa capacità.

\booksubsection{Scenario 5: ricerca immagini e riconoscimento oggetti — il tuo “motore di ricerca visivo”}

\textbf{Prestazioni reali}: vedi per strada un fiore che non conosci, apri la fotocamera e l'IA ti dice subito “è un ibisco”. Cerchi nell'album le foto scattate l'anno scorso al mare: cerca “mare” e l'IA trova automaticamente tutte le foto che lo contengono, senza sfogliare a mano.

\textbf{Principio alla base}: i modelli di visione artificiale (CV) trasformano il contenuto delle immagini in vettori (Embedding), da confrontare poi con la query testuale per similarità. Questa “ricerca multimodale” rappresenta un grande salto nella gestione dell'album del telefono.

\booksection{Ottimizzazione dei prompt: la “formula di comunicazione” per far ubbidire l'IA}

La qualità delle risposte dell'IA dipende in gran parte da come “poni la domanda”. Per la stessa domanda, modi diversi di chiedere producono risposte completamente diverse. Qui sotto presentiamo un approccio\textbf{molto pratico e facile da usare, basato su “Ruolo-Background-Compito-Vincoli”}, con molti esempi.

\begin{tcbinfo}{Perché serve un framework?}
Molti usano l'IA come usano Baidu — inserendo qualche parola chiave. Ma l'IA non è un motore di ricerca, è un modello linguistico che “legge, comprende e poi risponde”.\textbf{Dandole più contesto, dicendo chiaramente cosa vuoi e limitandola per evitare derive}, otterrai risultati precisi. Il framework serve a tradurre nella tua mente “come dev'essere una buona risposta” in istruzioni comprensibili all'IA.

\end{tcbinfo}

\booksubsection{Ruolo (Role): dì all'IA chi è}

Dando all'IA un'“identità” specifica, richiamerà le conoscenze professionali e lo stile espressivo del settore corrispondente.\textbf{Più specifico è il ruolo, più professionale sarà l'output}.

\booksubsection{Esempi pratici:}

\begin{itemize}
  \item "Sei un\textbf{tecnico riparazione telefoni Xiaomi con 10 anni di esperienza}." (problema tecnico)
  \item "Sei un\textbf{professore di fisica delle scuole medie bravo a spiegare in modo semplice concetti complessi}." (divulgazione scientifica)
  \item "Sei un\textbf{editor librario con buon gusto, molto colto e con uno stile tutto suo}." (scrittura)
  \item "Sei un\textbf{blogger di cucina che ogni giorno prepara la colazione per i lavoratori}." (vita quotidiana)
  \item "Sei un\textbf{medico di medicina tradizionale cinese paziente, gentile e che non deride mai i principianti}." (salute)
\end{itemize}

\textbf{Perché funziona?}Il ruolo definisce\textbf{l'ambito di conoscenza}, lo\textbf{stile espressivo}e la\textbf{profondità professionale}dell'IA. Un'IA senza ruolo è una “tuttologa che non eccelle in nulla”; con un ruolo, può “entrare” in quell'identità e pensare come quel professionista.

\booksubsection{Background: dì all'IA la tua situazione}

Spiega in anticipo “chi sei, qual è la tua situazione attuale, quali sono i tuoi vincoli”, e l'IA ti darà risposte\textbf{davvero adatte a te}, non generiche.

\booksubsection{Esempi pratici:}

\begin{itemize}
  \item "Il mio telefono è uno\textbf{Xiaomi 14, comprato sei mesi fa, ultimamente si scalda molto quando gioco}." (situazione dispositivo)
  \item "Sono\textbf{una signora in pensione di 60 anni, ci vedo poco e non so usare bene lo smartphone}." (profilo utente)
  \item "Sto\textbf{preparando l'esame di inglese di Livello 6 di giugno 2026, con una base media}." (background di studio)
  \item "Ho aperto\textbf{un piccolo bar in una città di terza fascia, un fatturato mensile di 30.000, con due dipendenti}." (contesto aziendale)
  \item "Ho\textbf{appena avuto un bambino e sono in casa per il puerperio, dormo troppo poco}." (fase di vita)
\end{itemize}

\textbf{Perché funziona?}Il background fa sapere all'IA\textbf{“per chi lavora”}e\textbf{“in quali condizioni”}. Per la stessa domanda, la spiegazione a uno studente o a un genitore è diversa; i consigli a un cittadino o a chi vive in paese sono diversi. Senza background, l'IA può solo dare la risposta del “minimo comun denominatore”, spesso poco incisiva.

\booksubsection{Compito (Task): dì chiaramente all'IA cosa fare}

In questa fase usa\textbf{un verbo all'inizio}(scrivi, elenca, confronta, traduci, spiega, riassumi...) e descrivi il compito in modo\textbf{concreto ed eseguibile}. Richieste vaghe portano solo a risposte vaghe.

\booksubsection{Esempi pratici:}

\begin{itemize}
  \item "\textbf{Elenca 5 autorizzazioni che dovrei revocare subito}e spiega perché." (azione chiara)
  \item "\textbf{Traduci in cinese la seguente email in inglese}, mantenendo un tono formale." (compito concreto)
  \item "\textbf{Confronta questi quattro telefoni intorno ai 3.000 yuan}e fornisci una tabella di pro e contro." (formato chiaro)
  \item "\textbf{Spiega con parole comprensibili a un bambino delle elementari 'cos'è il 5G'}." (stile chiaro)
  \item "\textbf{Scrivi tre testi per brevi video per il mio negozio di colazione}, rivolti rispettivamente a lavoratori, coppie e genitori con bambini." (segmentazione chiara)
\end{itemize}

\textbf{Perché funziona?}La sezione del compito è il “centro di comando” per le azioni dell'IA.\textbf{Più il verbo è chiaro e più le richieste sono concrete, più l'IA sa “cosa consegnare”}. Formule vaghe come “sistema un po'”, “fai una cosa”, “dai un'occhiata” sono il grande tabù per l'IA — interpreterà le tue intenzioni, spesso in modo errato.

\booksubsection{Vincoli (Constraints): dì all'IA cosa non fare}

I “vincoli” sembrano una “restrizione”, ma in realtà sono lo strumento per\textbf{portare l'IA a un output il più adatto possibile a te}. Le dimensioni comuni sono: numero di parole, formato, tono, cose da evitare, cose da includere.

\booksubsection{Esempi pratici:}

\begin{itemize}
  \item "\textbf{Non più di 200 parole, elenca per punti}." (lunghezza + formato)
  \item "\textbf{Tono cordiale e naturale, come una chiacchierata tra amici}, evita espressioni burocratiche come 'in primo luogo, in secondo luogo, in terzo luogo'." (tono)
  \item "\textbf{Non usare termini tecnici}, quando compaiono spiegali." (limitazione del lessico)
  \item "\textbf{Non più di 3 righe per paragrafo}, per agevolare la lettura sul telefono." (impaginazione)
  \item "\textbf{Deve includere: scenario d'uso, pro e contro, rispondenza al mio budget}." (elementi obbligatori)
  \item "\textbf{Non dare consigli di investimento specifici}, spiega solo i principi." (confini)
\end{itemize}

\textbf{Perché funziona?}Senza vincoli, l'IA tende a “improvvisare” — troppo accademica, troppo lunga o fuori tema. Vincoli chiari = “standard di consegna” chiari, allineano le aspettative e rendono i risultati controllabili.

\booksubsection{Valutazione della qualità del prompt: 5 caratteristiche di un buon prompt}

Scrivere bene i prompt non è un'operazione una tantum:\textbf{prima di inviarlo all'IA, verifica questi cinque punti}:

\begin{enumerate}
  \item \textbf{Concreto e non astratto}: non "scrivi un articolo", ma "scrivi una guida all'acquisto di telefoni di 800 parole, in 3 sezioni, rivolta a mamme alle prime armi".
  \item \textbf{Con background e non isolato}: dì all'IA "chi sono, perché lo chiedo", così potrà darti una risposta "personalizzata per te".
  \item \textbf{Con un formato di output chiaro}: tabella, elenco, paragrafo o JSON? Specificalo in anticipo.
  \item \textbf{Vincoli adeguati}: parole, stile, cose da evitare — sceglierne uno solo di questi tre porta già a un miglioramento evidente.
  \item \textbf{Verificabile dall'IA}: dopo aver chiesto, fai ripetere all'IA la tua richiesta; se non sa spiegarla, non era chiara.
\end{enumerate}

\booksubsection{Problemi comuni e soluzioni}

\begin{center}
\begin{tabularx}{\linewidth}{|X|X|}
\hline
\textbf{Sintomo} & \textbf{Soluzione} \tabularnewline
\hline
Risposta troppo generica, poco pratica & Aggiungi background (chi sei / in quale scenario) \tabularnewline
Risposta troppo accademica / troppo formale & Aggiungi vincoli di tono + una lista di espressioni da evitare \tabularnewline
Risposta fuori tema / non risponde alla domanda & Riformula il compito con un verbo e chiedi "rispondi solo a questo" \tabularnewline
Risposta troppo lunga, non finibile & Limita rigidamente le parole (massimo 200) \tabularnewline
L'IA finge con aria seria & Aggiungi "se non sei sicuro, scrivi 'non sono sicuro'" + verifica autonoma dei fatti importanti \tabularnewline
Prima versione insoddisfacente & Non riscrivere, aggiungi al prompt originale "un po' più..." \tabularnewline
\hline
\end{tabularx}
\end{center}

\booksubsection{Tre casi completi di scenario}

\begin{tcbinfo}{Scenario 1: assistente di vita (consigli per la colazione)}
\textbf{Prompt}:

Sei un blogger di colazione esperto di abbinamenti nutrizionali (\textbf{ruolo}). Sono una mamma lavoratrice, ho solo 15 minuti ogni mattina per preparare la colazione, a casa ho un bambino di 6 anni e una suocera di 60 (\textbf{background}).\textbf{Elenca 7 menu della colazione diversi per 7 giorni}(\textbf{compito}). Ogni colazione deve: (1) richiedere al massimo 15 minuti di preparazione (2) essere nutrizionalmente equilibrata (proteine + carboidrati + frutta e verdura) (3) essere adatta a tutti (4) con ingredienti reperibili in un supermercato normale (\textbf{vincoli}).

\textbf{Perché funziona}: il ruolo fissa il livello professionale, il background dà le caratteristiche comuni a tre categorie di persone, il compito definisce chiaramente il risultato, i 4 vincoli ne garantiscono l'eseguibilità.

\end{tcbinfo}

\begin{tcbinfo}{Scenario 2: supporto allo studio (preparazione all'esame di inglese di Livello 6)}
\textbf{Prompt}:

Sei un\textbf{professore di inglese di Livello 6 con molta esperienza}(ruolo). Sono uno studente universitario,\textbf{al Livello 4 ho preso 480, il mio obiettivo al Livello 6 è 500, ho ancora 3 mesi}(background).\textbf{Prepara per me un piano di studio di 12 settimane eseguibile}(compito). Requisiti: (1) ogni settimana con ripartizione chiara del tempo tra ascolto, parlato, lettura e scrittura (2) indica la direzione di miglioramento per ogni mese (3) consiglia 3 materiali di ripasso adatti (spiegando il motivo in cinese) (4) output in formato tabella (vincoli).

\textbf{Perché funziona}: il ruolo fissa il background professionale, il background fornisce i punti di ancoraggio di punteggio e tempo, il compito definisce chiaramente un “piano a fasi”, i vincoli richiedono “tabella” ed “eseguibile”.

\end{tcbinfo}

\begin{tcbinfo}{Scenario 3: efficienza sul lavoro (riassunto del verbale di riunione)}
\textbf{Prompt}:

Sei un\textbf{assistente amministrativo senior}(ruolo). Ho appena terminato una riunione di dipartimento di 1 ora,\textbf{sul piano di promozione marketing del Q2}(background).\textbf{Riassumi i seguenti appunti della riunione in un verbale strutturato}(compito). Requisiti: (1) dividi in “deliberazioni / azioni da fare / rischi e disaccordi” (2) per le azioni indica il responsabile e la scadenza (se non menzionati, scrivi “da confermare”) (3) massimo 600 parole in totale (4) output in formato markdown (vincoli).

\textbf{Perché funziona}: il ruolo fissa il punto di vista amministrativo, il background dà il tema della riunione, il compito definisce chiaramente l'output strutturato, i 4 vincoli rendono il risultato direttamente riutilizzabile.

\end{tcbinfo}

\textbf{Usi avanzati: rendi l'IA il tuo “assistente personale”}\par

Prova questi usi avanzati per sperimentare la vera forza dell'IA:

\begin{itemize}
  \item \textbf{Ottimizzazione multi-turno}: se la risposta dell'IA non ti soddisfa, puoi dire "è troppo tecnico, spiegalo di nuovo con parole comprensibili a un bambino delle elementari".
  \item \textbf{Chain of Thought (CoT)}: chiedi all'IA di "pensare passo dopo passo", in genere la qualità della risposta è superiore: "analizza passo dopo passo questo problema di matematica".
  \item \textbf{Collaborazione tra più IA}: usa A per la prima bozza, B per la rifinitura, C per la traduzione — combinarle dà risultati migliori.
  \item \textbf{Prompt di sistema}: preimposta nelle impostazioni dell'IA la tua identità, le preferenze, gli argomenti tabù, per avere risposte più in linea con le tue esigenze.
\end{itemize}

\booksection{Limiti e sfide dell'IA: mantenere un giudizio lucido}

Anche se potente, l'IA non è affatto “onnipotente”. Riconoscere lucidamente i suoi limiti è indispensabile per usarla meglio:

\begin{itemize}
  \item \textbf{Problema delle "allucinazioni"}: l'IA può dire fesserie con grande serietà (termine tecnico: Hallucination). Può inventare eventi storici inesistenti, citazioni errate, statistiche false. Le informazioni importanti vanno verificate manualmente.
  \item \textbf{Knowledge cutoff}: i dati di addestramento della maggior parte delle IA hanno una data di cutoff (ad es. inizio 2024); non sa cosa sia successo dopo.
  \item \textbf{Capacità di ragionamento limitata}: nel ragionamento multi-passaggio complesso, nelle dimostrazioni matematiche e nelle trappole logiche, l'IA sbaglia ancora facilmente.
  \item \textbf{Rischio per la privacy}: nei servizi IA cloud, le tue conversazioni possono essere usate per addestrare i modelli. Per le informazioni sensibili, preferisci l'IA on-device.
  \item \textbf{Pregiudizi e valori}: l'IA può riportare i pregiudizi presenti nei dati di addestramento; è necessario mantenere un pensiero indipendente.
\end{itemize}

\begin{tcbinfo}{Suggerimenti per l'uso dell'IA}
\begin{itemize}
  \item I contenuti generati dall'IA possono essere inaccurati (problema delle "allucinazioni"): verifica sempre le informazioni importanti.
  \item Non inserire nell'IA informazioni personali riservate (numero di carta d'identità, password, dati bancari, ecc.).
  \item Usa bene i prompt: più la descrizione è specifica e chiara, più alta è la qualità della risposta dell'IA.
  \item L'IA è uno strumento, non un'autorità: mantieni un pensiero e un giudizio indipendenti.
  \item Per decisioni cruciali (mediche, legali, finanziarie), consulta professionisti e non affidarti all'IA.
\end{itemize}

\end{tcbinfo}

\begin{tcbpractice}{Sfida: tre esercizi pratici con l'IA}
\begin{enumerate}
  \item \textbf{Sfida 1: istruzioni complesse}: dì all'assistente vocale: "Domani alle 7 ricordami di portare l'ombrello e sposta l'appuntamento delle 15 di oggi pomeriggio alle 16." Verifica se riesce a completare operazioni multi-step in una volta sola.
  \item \textbf{Sfida 2: fotoritocco IA}: scatta una foto turistica con passanti, usa la funzione "Eliminazione magica" o "Gomma IA" per rimuoverli. Osserva l'effetto di "ricostruzione" dello sfondo da parte dell'IA.
  \item \textbf{Sfida 3: ottimizzazione dei prompt}: prima chiedi all'IA con un prompt semplice "scrivi una guida di viaggio", poi usa il framework CRISPE per scrivere un prompt dettagliato. Confronta la qualità delle due risposte.
\end{enumerate}

\end{tcbpractice}

\begin{tcbinfo}{Caso 2: flusso completo di migrazione dati al cambio telefono}
\begin{itemize}
  \item \textbf{Scenario}: iPhone 13 Xiaomi 15, cambio cross-brand.
  \item \textbf{Passaggi}: (1) backup su iCloud/WeChat integrato del vecchio telefono (2) logout dell'Apple ID + disattivazione di "Trova" (3) sul nuovo telefono scansione con l'app "Xiaomi Switch" (4) selezione degli elementi da migrare (5) migrazione separata di WeChat con "Migrazione cronologia chat" (6) verifica dei dati chiave (contatti/album/WeChat) (7) cancellazione crittografata del vecchio telefono.
  \item \textbf{Risultato atteso}: nessuna perdita di dati, il vecchio telefono può essere rivenduto in sicurezza.
\end{itemize}

\end{tcbinfo}
